\documentclass[12pt,a4paper]{article}
\usepackage[utf8]{inputenc}
\usepackage[T1]{fontenc}
\usepackage[spanish]{babel}
\usepackage{amsmath, amssymb, amsfonts}
\usepackage{graphicx}
\usepackage{geometry}
\usepackage{booktabs}
\usepackage{enumitem}

\geometry{margin=2.5cm}

\begin{document}

% ----------------------------------------------------------------------
% ENCABEZADO INSTITUCIONAL
% ----------------------------------------------------------------------
\begin{center}
    \textbf{\Large UNIVERSIDAD AUTÓNOMA GABRIEL RENÉ MORENO} \\[0.2cm]
    \textbf{\large FACULTAD DE CIENCIAS CONTABLES} \\[0.2cm]
    \textbf{\large CARRERA DE CONTADURÍA PÚBLICA} \\[0.5cm]
    \textbf{\large ESTADÍSTICA II MAT-260} \\[1cm]
\end{center}

\tableofcontents
\vspace{1cm}
\hrule
\vspace{1cm}

% ----------------------------------------------------------------------
% UNIDAD 1
% ----------------------------------------------------------------------
\section{Unidad 1: Análisis Combinatorio}

\subsection{Introducción}
El \textbf{análisis combinatorio} es la rama de las matemáticas que trata los diversos tipos de arreglos o selecciones que se pueden formar con los elementos de un conjunto. 
El estudio del análisis combinatorio sirve como base para comprender y resolver problemas sobre probabilidades, ya que nos permite contar de manera eficiente los resultados posibles de un experimento aleatorio sin necesidad de enumerarlos uno por uno.

\subsection{Conteo}
El proceso de conteo simplifica la cuantificación de eventos. Para ello, nos apoyamos en dos principios fundamentales:

\textbf{Ley de la multiplicación:} Si una elección consta de $k$ pasos, donde el primero puede realizarse de $n_{1}$ maneras, el segundo de $n_{2}$ maneras, y así sucesivamente, el número total de formas en que toda la situación puede resolverse es el producto de todas estas formas.
\[
\text{Total de formas} = n_{1} \times n_{2} \times \dots \times n_{k}
\]
Esta ley aplica para eventos que no son mutuamente excluyentes (pueden ocurrir en forma conjunta de manera secuencial).

\textbf{Ley de la adición:} Si un evento $A$ puede ocurrir de $n$ formas diferentes y un evento $B$ de $m$ formas diferentes, cumpliendo la condición de que los eventos $A$ y $B$ son mutuamente excluyentes (la ocurrencia de uno impide la del otro), el número total de formas es la suma de sus formas individuales.
\[
\text{Total de formas} = n + m
\]

\vspace{0.5cm}
\textbf{Ejemplo resuelto 1: Ley de la multiplicación} \\
Un estudiante de Contaduría Pública desea armar un conjunto de ropa formal para una presentación. Tiene 4 camisas distintas, 3 pantalones de vestir y 2 pares de zapatos. ¿De cuántas formas diferentes puede vestirse?

\textbf{Solución paso a paso:}
\begin{enumerate}[noitemsep]
    \item Identificamos el número de opciones para el primer paso (camisas): $n_{1} = 4$.
    \item Identificamos el número de opciones para el segundo paso (pantalones): $n_{2} = 3$.
    \item Identificamos el número de opciones para el tercer paso (zapatos): $n_{3} = 2$.
    \item Dado que la elección de una prenda no excluye a la otra, aplicamos la ley de la multiplicación:
    \[
    \text{Formas totales} = 4 \times 3 \times 2 = 24
    \]
\end{enumerate}
El estudiante puede armar 24 conjuntos de ropa formal diferentes.

\subsection{Permutaciones}
Una \textbf{permutación} es un arreglo ordenado de elementos. El número de permutaciones de $n$ objetos es el número de formas en las que pueden acomodarse esos objetos considerando estrictamente el orden de colocación.

\textbf{Permutaciones de $n$ objetos (Factorial):} El número total de arreglos posibles de $n$ objetos tomados todos a la vez es $n$ factorial, denotado por $P_{n} = n!$.
\[
P_{n} = n! = n \times (n-1) \times (n-2) \times \dots \times 3 \times 2 \times 1
\]
\textbf{Propiedades del factorial:} $0! = 1$ y $1! = 1$.

\textbf{Permutaciones de $n$ objetos tomados de $r$ a la vez:} Cuando nos interesa el número de arreglos de un subgrupo de tamaño $r$ extraído de un conjunto mayor $n$, utilizamos la fórmula:
\[
_{n}P_{r} = \frac{n!}{(n-r)!}
\]

\textbf{Permutaciones con repetición:} Si de $n$ objetos, $r_{1}$ son idénticos, $r_{2}$ son idénticos y así sucesivamente, el número de permutaciones distinguibles es:
\[
P_{\text{repetición}} = \frac{n!}{r_{1}! \times r_{2}! \times \dots \times r_{k}!}
\]

\textbf{Permutaciones circulares:} El número de formas en que $n$ objetos pueden acomodarse en un círculo (donde el punto de inicio es relativo) se calcula como:
\[
P_{\text{circular}} = (n-1)!
\]

\vspace{0.5cm}
\textbf{Ejemplo resuelto 2: Permutaciones de $n$ tomados de $r$} \\
Una firma auditora debe asignar a 3 de sus 8 mejores auditores a tres sucursales distintas: una en La Paz, una en Cochabamba y otra en Santa Cruz. ¿De cuántas formas distintas se puede hacer la asignación?

\textbf{Solución paso a paso:}
\begin{enumerate}[noitemsep]
    \item Identificamos que el orden importa, ya que asignar al auditor A a La Paz y al auditor B a Santa Cruz es un arreglo diferente a asignar a B a La Paz y a A a Santa Cruz.
    \item Identificamos la población total $n = 8$ y los puestos a ocupar $r = 3$.
    \item Aplicamos la fórmula de permutaciones sin repetición:
    \[
    _{8}P_{3} = \frac{8!}{(8-3)!} = \frac{8!}{5!}
    \]
    \item Desarrollamos el factorial de 8 hasta el 5 para poder simplificar:
    \[
    _{8}P_{3} = \frac{8 \times 7 \times 6 \times 5!}{5!} = 8 \times 7 \times 6 = 336
    \]
\end{enumerate}
Existen 336 formas diferentes de asignar a los tres auditores a las sucursales correspondientes.

\subsection{Combinaciones}
Una \textbf{combinación} es una selección o agrupación de objetos sin tener en cuenta el orden en que son seleccionados. A diferencia de las permutaciones, en las combinaciones solo nos interesa qué elementos componen el grupo, no cómo están acomodados.

\textbf{Combinaciones de $n$ objetos tomados de $r$ a la vez:} El número de combinaciones posibles se obtiene dividiendo las permutaciones entre las formas de ordenar el subgrupo $r$:
\[
_{n}C_{r} = \frac{n!}{r! (n-r)!}
\]

\vspace{0.5cm}
\textbf{Ejemplo resuelto 3: Combinaciones} \\
Para revisar y auditar los balances de una empresa, se requiere formar un equipo de 4 analistas. Si el departamento cuenta con 10 analistas disponibles, ¿de cuántas formas se puede seleccionar el equipo de trabajo?

\textbf{Solución paso a paso:}
\begin{enumerate}[noitemsep]
    \item Observamos que el orden de selección no es relevante (el equipo formado por Ana, Luis, Carlos y María es exactamente el mismo sin importar quién fue nombrado primero). Usamos combinaciones.
    \item Determinamos la población $n = 10$ y el tamaño del subgrupo $r = 4$.
    \item Reemplazamos los valores en la fórmula de combinaciones:
    \[
    _{10}C_{4} = \frac{10!}{4! (10-4)!} = \frac{10!}{4! \times 6!}
    \]
    \item Expandimos el numerador hasta el 6 factorial para facilitar la simplificación:
    \[
    _{10}C_{4} = \frac{10 \times 9 \times 8 \times 7 \times 6!}{4 \times 3 \times 2 \times 1 \times 6!}
    \]
    \item Simplificamos el cálculo aritmético multiplicando el denominador ($4 \times 3 \times 2 \times 1 = 24$) y operando:
    \[
    _{10}C_{4} = \frac{5040}{24} = 210
    \]
\end{enumerate}
El gerente puede formar 210 equipos de analistas diferentes.

\subsection{Diferencia entre Permutación y Combinación}
La principal diferencia y la clave para resolver cualquier ejercicio de la unidad es \textbf{la importancia del orden}. 
A cada grupo ordenado de elementos se le llama permutación, mientras que a un grupo no ordenado se le llama combinación. 

Para explicarlo de manera evidente, consideremos un conjunto de letras $\{A, B, C\}$. 
Si deseamos seleccionar grupos de 3 letras evaluando cada caso, tenemos:
\begin{itemize}
    \item \textbf{Permutación (el orden importa):} Obtendremos 6 arreglos diferentes: $(A,B,C)$, $(A,C,B)$, $(B,A,C)$, $(B,C,A)$, $(C,A,B)$, y $(C,B,A)$. Cada uno representa una permutación diferente por su orden de colocación.
    \item \textbf{Combinación (el orden no importa):} Esas mismas letras representan exactamente un mismo grupo en su naturaleza elemental. Por tanto, las 6 permutaciones representan $1$ sola combinación.
\end{itemize}

% ----------------------------------------------------------------------
% EJERCICIOS PROPUESTOS
% ----------------------------------------------------------------------
\subsection{Ejercicios propuestos}
\begin{enumerate}
    \item Una clave de acceso al aula virtual Moodle debe estar formada por 3 vocales distintas seguidas de 2 dígitos pares distintos. Utilizando la ley de la multiplicación, ¿cuántas contraseñas diferentes pueden crearse?
    
    \item En una mesa redonda de la Facultad de Ciencias Contables deben sentarse el decano, el vicedecano y 4 directores de carrera. ¿De cuántas formas distintas pueden acomodarse considerando la permutación circular?
    
    \item De un grupo de 15 estudiantes destacados de Estadística II, el docente debe seleccionar a un comité de 4 personas para coordinar un seminario. ¿De cuántas maneras diferentes puede seleccionar al grupo si no hay jerarquía entre ellos?
    
    \item Con las letras de la palabra "PROBABILIDAD", ¿cuántas permutaciones diferentes se pueden formar considerando que existen letras repetidas?
    
    \item Se desea conformar una comisión de presupuesto compuesta por 3 profesores titulares y 2 profesores interinos. Si hay 8 profesores titulares y 5 profesores interinos disponibles, ¿de cuántas maneras se puede establecer dicha comisión?
\end{enumerate}

\end{document}